\documentclass{article}
\usepackage{amsmath, amssymb, amsthm}
\usepackage{hyperref}
\usepackage{geometry}
\geometry{margin=1in}

\title{Erd\H{o}s Problem \#393}
\date{Last updated: 2025-08-31}
\author{Source: erdosproblems.com}

\begin{document}
\maketitle

\noindent\textbf{Status:} OPEN \quad \textbf{No prize}

\noindent\textbf{Tags:} number theory, factorials

\medskip

\noindent\textbf{Problem Statement:}

Let $f(n)$ denote the minimal $m\geq 1$ such that\[n! = a_1\cdots a_t\]with $a_1&#60;\cdots &#60;a_t=a_1+m$. What is the behaviour of $f(n)$?




Erd\H{o}s and Graham write that they do not even know whether $f(n)=1$ infinitely often (i.e. whether a factorial is the product of two consecutive integers infinitely often). Let $F_m(N)$ count the number of $n\leq N$ such that $f(n)=m$. Berend and Osgood \cite{BeOs92} proved that, for each fixed $m$, $F_m(N)=o(N)$. Bui, Pratt, and Zaharescu \cite{BPZ23} have shown that\[F_m(N)\ll_m N^{33/34}.\]A result of Luca \cite{Lu02} implies that $f(n)\to \infty$ as $n\to \infty$, conditional on the ABC conjecture .




References

[BPZ23] Bui, Hung M. and Pratt, Kyle and Zaharescu, Alexandru, Power savings for counting solutions to polynomial-factorial equations . Adv. Math. (2023), Paper No. 109021, 32.

[BeOs92] Berend, Daniel and Osgood, Charles F., On the equation {$P(x)=n!$} and a question of {E}rd\H{o}s . J. Number Theory (1992), 189--193.

[Lu02] Luca, Florian, The {D}iophantine equation {$P(x)=n!$} and a result of {M}. {O}verholt . Glas. Mat. Ser. III (2002), 269--273.

\noindent\textbf{Related OEIS sequences:} A388302


\bigskip
\noindent\small{Source: \url{https://www.erdosproblems.com/393}}
\end{document}
